\subsection{Adquisición de datos e instancias}

La evaluación del modelo se apoya en instancias estándar de la biblioteca
\textit{Capacitated Vehicle Routing Problem Library} (CVRPLIB), repositorio ampliamente
utilizado para contrastar algoritmos exactos sobre un conjunto común de casos de prueba
y soluciones de referencia \parencite{cvrplib}. El uso de esta biblioteca permite
trabajar sobre instancias estructuradas y comparables, lo que resulta especialmente
importante en un problema $\mathcal{NP}$-hard como el CVRP, donde la dificultad
computacional crece rápidamente con el número de nodos y con la rigidez de las
restricciones de capacidad.

El flujo de adquisición de datos considera dos etapas. En primer lugar, las
instancias del Set A y Set E fueron obtenidas desde el repositorio oficial de CVRPLIB
y almacenadas localmente junto con sus soluciones de referencia (archivos
\texttt{.vrp} y \texttt{.sol}), lo que asegura trazabilidad y reproducibilidad
en la construcción del conjunto experimental. En segundo lugar, cada archivo
\texttt{.vrp} es procesado para extraer la información relevante del problema:
número de nodos, capacidad vehicular, coordenadas cartesianas del depósito y
de los clientes, y demanda asociada a cada nodo de atención. El número de
vehículos $m$ se obtiene directamente del identificador de la instancia
(componente $k$ del nombre), y el valor óptimo de referencia se lee desde el
archivo \texttt{.sol} correspondiente. A partir de estos datos se construye la
estructura del grafo completo no dirigido y se calcula la matriz de costos
$c_{ij}$ empleada por el modelo.

La evaluación considera cinco instancias de los Sets E y A de CVRPLIB en orden
creciente de tamaño: \texttt{E-n13-k4} ($n=13$, $m=4$), \texttt{E-n22-k4}
($n=22$, $m=4$), \texttt{E-n23-k3} ($n=23$, $m=3$), \texttt{A-n32-k5} ($n=32$,
$m=5$) y \texttt{A-n33-k5} ($n=33$, $m=5$). Las tres primeras, del Set E 
\parencite{christofides1981exact}, son de tamaño pequeño y permiten validar la
correctitud del procedimiento; \texttt{A-n32-k5}, caso clásico de referencia
para métodos exactos del Set A \parencite{augerat1995}, constituye la
instancia objetivo del estudio, y \texttt{A-n33-k5} explora el comportamiento
del método un escalón por sobre ella. Esta progresión de tamaños permite
observar el crecimiento del esfuerzo computacional y delimitar empíricamente el
rango de instancias resolubles a optimalidad dentro de un presupuesto de tiempo
dado.

Siguiendo la convención \texttt{EUC\_2D} de TSPLIB utilizada por CVRPLIB, el
costo de cada arista se define como la distancia euclidiana redondeada al
entero más cercano, $c_{ij} = \left[ \sqrt{(x_i-x_j)^2 + (y_i-y_j)^2} \right]$,
lo que garantiza la comparabilidad directa con los valores óptimos reportados
en la biblioteca.

Algunas instancias del Set E declaran sus distancias de forma explícita
(\texttt{EDGE\_WEIGHT\_TYPE: EXPLICIT}, con la matriz triangular inferior en el
propio archivo); en estos casos, la matriz de costos se lee directamente del
archivo en lugar de calcularse desde coordenadas. El procedimiento de carga
soporta ambos formatos de manera transparente para el modelo; cuando la
instancia carece de coordenadas, se omite únicamente la visualización de rutas,
sin efecto sobre la resolución.

\subsection{Formulación del modelo matemático}

Sea $G=(V,E)$ un grafo completo no dirigido, donde $V=\{1,2,\dots,n\}$ es el conjunto
de nodos, con el nodo $1$ representando el depósito, y $E$ es el conjunto de aristas no
dirigidas. Cada cliente $i\in V\setminus\{1\}$ posee una demanda $d_i \ge 0$, y cada
arista $\{i,j\}\in E$ tiene un costo de recorrido $c_{ij}=c_{ji}$. Se considera una
flota homogénea de $m$ vehículos, todos con capacidad $D$.

\subsubsection{Variables de decisión}

Siguiendo la formulación simétrica propuesta por \textcite{laporte1983branch}, se
define la variable de decisión $x_{ij}$ como el número de veces que una solución
utiliza la arista no dirigida $\{i,j\}$, para todo $i<j$.

\subsubsection{Función objetivo}

La función objetivo minimiza la distancia total recorrida por la flota:
\begin{equation}
    \minimo \; z = \sum_{\{i,j\} \in E} c_{ij}\, x_{ij}.
    \label{eq:metodo_objetivo}
\end{equation}

\subsubsection{Restricciones}

La estructura de las rutas queda determinada por restricciones de grado. Para el
depósito, la suma de incidencias debe ser igual al doble del número de vehículos
utilizados:
\begin{equation}
    \sum_{j=2}^{n} x_{1j} = 2m.
    \label{eq:metodo_deposito}
\end{equation}

Para cada cliente $k=2,\dots,n$, la solución debe asegurar grado dos, lo que equivale
a imponer que cada nodo de demanda sea visitado exactamente una vez dentro de alguna
ruta:
\begin{equation}
    \sum_{i<k} x_{ik} + \sum_{j>k} x_{kj} = 2,
    \quad k=2,\dots,n.
    \label{eq:metodo_clientes}
\end{equation}

Las restricciones \eqref{eq:metodo_deposito} y \eqref{eq:metodo_clientes} son
necesarias pero no suficientes para garantizar factibilidad global, ya que admiten
soluciones que contienen ciclos desconectados del depósito. Para descartar estas
configuraciones e incorporar simultáneamente la capacidad vehicular, se emplean
restricciones de eliminación de subtours y capacidad sobre todo subconjunto de
clientes $S \subseteq V \setminus \{1\}$ con $|S| \ge 2$. La forma original propuesta
por \textcite{laporte1983branch} expresa estas restricciones como cortes de capacidad:
\begin{equation}
    \sum_{i \in S} \sum_{j \notin S} x_{ij} \;\ge\; 2 \left\lceil
        \frac{\sum_{i \in S} d_i}{D}
    \right\rceil,
    \quad S \subseteq V \setminus \{1\},\; |S| \ge 2.
    \label{eq:metodo_corte}
\end{equation}

En presencia de las restricciones de grado, la desigualdad
\eqref{eq:metodo_corte} es equivalente a la \textit{rounded capacity inequality} (RCI),
empleada habitualmente en implementaciones modernas
\parencite{lysgaard2004branch}:
\begin{equation}
    \sum_{i,j \in S\,;\, i<j} x_{ij}
    \;\le\; |S| - \left\lceil
        \frac{\sum_{i \in S} d_i}{D}
    \right\rceil,
    \quad S \subseteq V \setminus \{1\},\; |S| \ge 2.
    \label{eq:metodo_rci}
\end{equation}

La desigualdad \eqref{eq:metodo_rci} expresa que, si la demanda agregada del
subconjunto $S$ requiere al menos $\left\lceil \sum_{i \in S} d_i / D \right\rceil$
vehículos, entonces no es posible que todas las aristas de $S$ queden contenidas en un
único subtour aislado. En consecuencia, esta familia de restricciones cumple una
doble función: descartar subtours ilegales y asegurar que la partición de clientes en
rutas respete la capacidad de carga disponible.

\subsubsection{Dominio de las variables}

El dominio de las variables refleja la estructura simétrica del modelo
\parencite{laporte1983branch}:
\begin{equation}
    x_{ij} \in \{0, 1\}, \quad \forall\, i, j \in V \setminus \{1\},\; i < j,
    \label{eq:metodo_dominio_clientes}
\end{equation}
\begin{equation}
    x_{1j} \in \{0, 1, 2\}, \quad \forall\, j \in V \setminus \{1\}.
    \label{eq:metodo_dominio_deposito}
\end{equation}

Las aristas entre clientes son binarias porque cada par de nodos puede recorrerse a lo
más una vez en una solución factible. En cambio, las aristas incidentes al depósito
admiten el valor $2$, lo que corresponde al caso particular en que una ruta atiende a
un único cliente y, por lo tanto, recorre la arista depósito-cliente dos veces (ida y
vuelta).

Desde el punto de vista computacional, el elemento más crítico del modelo no es
únicamente la integralidad, sino la naturaleza exponencial de la familia
\eqref{eq:metodo_rci}, lo que hace inviable incorporarla exhaustivamente desde el
inicio para instancias de tamaño no trivial.

\subsection{Implementación computacional}

La implementación se desarrolló en \textbf{Python}, utilizando la biblioteca
\textbf{PuLP} como interfaz de modelamiento para programación lineal entera
mixta \parencite{pulp2011}. Las pruebas iniciales y la validación del
procedimiento se realizaron en \textbf{Google Colab} (Python 3.10), mientras
que la experimentación definitiva y las ejecuciones finales reportadas en este
informe se llevaron a cabo en un entorno local (Python 3.12). Esta elección
permite construir el modelo de manera declarativa, mantener una representación
transparente de las variables y restricciones, e integrar con facilidad un
esquema iterativo de fortalecimiento del modelo; el código completo se encuentra disponible en un
repositorio público\footnote{\url{https://github.com/OrejiNet/CVRP_BNB}}.

La experimentación definitiva se ejecutó en un equipo de escritorio con
procesador Intel Core i9-11900K (3{,}50~GHz, 8 núcleos), 32~GB de memoria RAM
y sistema operativo Windows 11.


La resolución del problema entero se delega al solver \textbf{COIN-OR CBC},
invocado desde PuLP \parencite{pulp2011}. CBC implementa internamente el algoritmo de
\textit{Branch and Bound} sobre la formulación entera, complementado con planos de
corte genéricos. Sobre esta capa, el procedimiento implementado replica la lógica
exacta del enfoque clásico adoptado \parencite{laporte1983branch}: se parte con una relajación del
modelo que considera únicamente la función objetivo \eqref{eq:metodo_objetivo} y las
restricciones de grado \eqref{eq:metodo_deposito}--\eqref{eq:metodo_clientes}. Tras
resolver el problema, se inspecciona la solución obtenida para determinar si la
estructura propuesta corresponde efectivamente a un conjunto de rutas factibles.

Para esta etapa de verificación se utiliza \textbf{NetworkX}
\parencite{hagberg2008exploring}, biblioteca empleada para representar la solución
como un grafo y detectar componentes conexas inducidas por las aristas activas. Esta
representación permite identificar de forma directa subtours desconectados del
depósito o componentes cuya demanda total excede la capacidad vehicular disponible. En
cada iteración, para cada subconjunto $S$ que incumpla la restricción de capacidad o
que induzca un subtour ilegal, se agrega al modelo una nueva desigualdad de la
forma \eqref{eq:metodo_rci}, y el problema se resuelve nuevamente. El ciclo continúa
hasta que ningún subconjunto requiera nuevas restricciones, momento en el cual la
solución resultante satisface simultáneamente las condiciones de conectividad y
capacidad.

Cada resolución del modelo entero se ejecuta con un límite de 60 segundos; si
el solver agota este presupuesto sin certificar la optimalidad del subproblema,
retorna la mejor solución entera encontrada, lo que se refleja en el análisis
de resultados del Capítulo~4.


\begin{figure}[H]
    \centering
    \shorthandoff{<>}
    \resizebox{0.84\textwidth}{!}{%
    \begin{tikzpicture}[node distance=0.75cm]
        \tikzstyle{proceso}=[rectangle, rounded corners, draw=black, thick, fill=gray!15, text width=4.8cm, minimum height=0.9cm, align=center]
        \tikzstyle{datos}=[trapezium, trapezium left angle=70, trapezium right angle=110, draw=black, thick, fill=blue!15, text width=4.8cm, minimum height=0.9cm, align=center]
        \tikzstyle{decision}=[diamond, draw=black, thick, fill=yellow!25, text width=3.2cm, aspect=2, align=center, inner sep=1pt]
        \tikzstyle{termino}=[ellipse, draw=black, thick, fill=gray!20, minimum width=2.4cm, minimum height=0.8cm, align=center]

        \node (inicio) [termino] {Inicio};
        \node (datos) [datos, below=of inicio] {Carga de la instancia \texttt{.vrp} y de su solución de referencia \texttt{.sol} desde repositorio local (Set A,Set B, CVRPLIB)};
        \node (parsing) [proceso, below=of datos] {Pre-procesamiento: \textit{parsing} del archivo \texttt{.vrp} y extracción de coordenadas, demandas, capacidad $D$ y número de vehículos $m$};
        \node (distancias) [proceso, below=of parsing] {Cálculo de la matriz de distancias euclidianas redondeadas (\texttt{EUC\_2D}) entre todos los pares de nodos};
        \node (modelo) [proceso, below=of distancias] {Configuración del modelo en PuLP: función objetivo y restricciones de grado iniciales};
        \node (resolver) [proceso, below=0.9cm of modelo] {Resolución del modelo entero con el solver COIN-OR CBC};
        \node (decision) [decision, below=1.0cm of resolver] {Subtours ilegales o exceso de capacidad detectados con NetworkX};
        \node (cortes) [proceso, right=3.5cm of decision] {Generación e incorporación de cortes dinámicos RCI (Ec.~\ref{eq:metodo_rci})};
        \node (nota) [proceso, below=0.8cm of cortes, text width=4.6cm, fill=blue!8] {\textbf{Ec.~(\ref{eq:metodo_rci})}\\[2pt] $\sum_{i,j\in S;\, i<j} x_{ij} \le |S| - \left\lceil \frac{\sum_{i\in S} d_i}{D} \right\rceil$};
        \node (salida) [proceso, below=1.1cm of decision] {Extracción de rutas finales y verificación de factibilidad};
        \node (visualizacion) [proceso, below=of salida] {Visualización de trayectorias con Matplotlib};
        \node (fin) [termino, below=of visualizacion] {Fin del proceso};

        \draw[->, thick] (inicio) -- (datos);
        \draw[->, thick] (datos) -- (parsing);
        \draw[->, thick] (parsing) -- (distancias);
        \draw[->, thick] (distancias) -- (modelo);
        \draw[->, thick] (modelo) -- (resolver);
        \draw[->, thick] (resolver) -- (decision);
        \draw[->, thick] (decision) -- node[above] {Sí} (cortes);
        \draw[->, thick] (decision) -- node[right] {No} (salida);
        \draw[->, thick] (cortes.north) |- node[pos=0.25, above] {reoptimizar} (resolver.east);
        \draw[->, thick] (cortes) -- (nota);
        \draw[->, thick] (salida) -- (visualizacion);
        \draw[->, thick] (visualizacion) -- (fin);
    \end{tikzpicture}%
    }
    \shorthandon{<>}
    \caption{Flujo de trabajo computacional para la resolución exacta del CVRP mediante generación iterativa de desigualdades de capacidad redondeada (RCI). Fuente: elaboración propia.}
    \label{fig:flujo_cvrp}
\end{figure}


\subsection{Procedimiento iterativo de generación de cortes}

El esquema computacional adoptado puede resumirse en los siguientes pasos:
\begin{enumerate}
    \item Cargar y parsear la instancia \texttt{.vrp} desde el repositorio local, junto con el valor óptimo de referencia del archivo \texttt{.sol}.
    \item Construir el grafo completo y calcular la matriz de distancias $c_{ij}$ a
          partir de las coordenadas euclidianas de los nodos.
    \item Formular un modelo inicial con la función objetivo
          \eqref{eq:metodo_objetivo} y las restricciones de grado
          \eqref{eq:metodo_deposito}--\eqref{eq:metodo_clientes}.
    \item Resolver el modelo entero con CBC y extraer las aristas activas de la
          solución.
    \item Analizar la solución mediante NetworkX para identificar componentes conexas
          distintas del depósito o rutas cuya carga acumulada exceda la capacidad $D$.
    \item Agregar dinámicamente las restricciones RCI correspondientes a los
          subconjuntos identificados y reoptimizar el modelo.
    \item Repetir el procedimiento hasta obtener una solución entera sin subtours
          ilegales y con factibilidad de carga en cada ruta.
\end{enumerate}

Este enfoque evita incorporar de forma anticipada un número exponencial de
restricciones, reduciendo el tamaño inicial del modelo y permitiendo que el
fortalecimiento ocurra únicamente cuando la estructura de la solución lo exige. Desde
el punto de vista de optimización exacta, esta estrategia puede interpretarse como un
esquema de generación iterativa de cortes sobre una formulación entera, estrechamente
vinculado al espíritu del método de \textit{Branch and Bound} propuesto por
\textcite{laporte1983branch}. La principal dificultad esperada radica en que, a medida
que crece el tamaño de la instancia, aumenta el número potencial de subconjuntos
candidatos a incumplir las restricciones de capacidad, lo que refleja directamente la
complejidad combinatoria inherente del CVRP.